A hyperbolic plane requires one plus and one minus among the representative's directions, and a definite signature has none. Otherwise Witt's theorem (\claimref{std-algebra.witt-extension}) extends any isometry between nondegenerate subspaces to the whole space; two hyperbolic planes are isometric, so the isometry group is transitive on them, and composing with the isometry of a plane that exchanges its null lines gives transitivity on null lines. In $(1,3)$ the rotation $e_{1} \mapsto e_{2}$, $e_{2} \mapsto -e_{1}$ fixes $e_{0}$ and its orientation and carries $\mathrm{span}(e_{0}, e_{1})$ onto $\mathrm{span}(e_{0}, e_{2})$ and $\R(e_{0} + e_{1})$ onto $\R(e_{0} + e_{2})$: the time orientation distinguishes a time direction, not a plane or a null line, the candidate flags through a time axis forming the celestial sphere. In $(2,2)$ the totally isotropic planes form a one-parameter family per component and the hyperbolic planes a four-parameter family, likewise permuted transitively. The script checks the witnesses and the Jacobian rank exactly. The obstruction is the same in every nondefinite residue, so it cannot sterilise $r = 0$ while leaving $r = 6$ causal.
